\ExerciseChapter{Logistic回归}

\label{chapter:2}

\par\addvspace{8pt}\noindent\begin{minipage}{\linewidth}

\section{作业题}

\subsection{数据分析题}

\Question{H2-1}{食管癌与饮酒：分层频数资料}

表11 按年龄分层后食管癌与饮酒关系。每日饮酒量单位为 g/天。\textbf{共7分。}

\begin{center}\small\setlength{\tabcolsep}{4pt}\renewcommand{\arraystretch}{1.22}\begin{tabular}{@{}lrrrr@{}}
\toprule
年龄组 & 0--79 病例 & 0--79 对照 & 80+ 病例 & 80+ 对照 \\
\midrule

25--44 & 5 & 35 & 5 & 270 \\
45--54 & 25 & 29 & 21 & 138 \\
55--64 & 42 & 27 & 34 & 139 \\
65+ & 24 & 18 & 44 & 119 \\
合计 & 96 & 109 & 104 & 666 \\
\bottomrule\end{tabular}\end{center}

\begin{enumerate}
\def\labelenumi{\arabic{enumi}.}
\tightlist
\item
  构建 Logistic 回归模型。（1分）
\item
  解释回归系数的含义。（2分）
\item
  如何检验总体回归系数是否为零？（2分）
\item
  如何评价模型的好坏？（2分）
\end{enumerate}


\worklines{3}

\end{minipage}\par\vspace{5pt}

\par\addvspace{8pt}\noindent\begin{minipage}{\linewidth}

\subsection{文献分析题}

\Question{H2-2}{睡眠质量与认知功能：文献分析}

阅读 \emph{Sleep quality and cognitive impairment in older Chinese: Guangzhou Biobank Cohort Study}（本包``数据与代码''附原文）。\textbf{共3分。}

\begin{enumerate}
\def\labelenumi{\arabic{enumi}.}
\tightlist
\item
  请阐述本研究的研究设计。（0.5分）
\item
  本研究的结局变量是如何测量及定义的？（1分）
\item
  请阐述 Table 2 中睡眠相关因素与记忆障碍的相关关系。（1.5分）
\end{enumerate}

Table 2（记忆障碍部分，保留全部因素；下列为 OR及95\% CI）：

\begin{center}\small\setlength{\tabcolsep}{4pt}\renewcommand{\arraystretch}{1.22}\begin{tabular}{@{}lll@{}}
\toprule
因素及比较 & Model 1 & Model 2 \\
\midrule

主观睡眠差 vs 好 & 1.09 (0.93--1.29) & 1.01 (0.85--1.19) \\
入睡31--60 vs \(\le30\) min & 1.09 (0.93--1.28) & 1.04 (0.89--1.22) \\
入睡\(>60\) vs \(\le30\) min & 1.23 (1.00--1.51)* & 1.13 (0.92--1.40) \\
入睡时间趋势 \(P\) & 0.04 & 0.24 \\
睡眠\(<5\) vs 7--8 h & 1.16 (0.90--1.48) & 1.12 (0.87--1.44) \\
睡眠5--6 vs 7--8 h & 0.99 (0.87--1.13) & 0.99 (0.87--1.13) \\
睡眠\(>8\) vs 7--8 h & 0.88 (0.63--1.22) & 0.90 (0.65--1.26) \\
睡眠时长趋势 \(P\) & 0.65 & 0.70 \\
效率75--85\% vs \(\ge85\%\) & 1.08 (0.92--1.27) & 1.09 (0.93--1.28) \\
效率65--75\% vs \(\ge85\%\) & 1.04 (0.84--1.29) & 1.02 (0.82--1.26) \\
效率\(<65\%\) vs \(\ge85\%\) & 1.26 (1.03--1.54)* & 1.23 (1.00--1.50)* \\
睡眠效率趋势 \(P\) & 0.04 & 0.08 \\
睡眠干扰每周至少1次 & 1.11 (0.77--1.59) & 0.89 (0.61--1.29) \\
使用睡眠药物 vs 不使用 & 1.02 (0.76--1.37) & 0.99 (0.73--1.33) \\
日间障碍每周至少1次 & 1.26 (1.02--1.56)* & 1.05 (0.85--1.31) \\
PSQI整体睡眠差 vs 好 & 1.10 (0.97--1.26) & 1.03 (0.90--1.18) \\
\bottomrule\end{tabular}\end{center}

睡眠干扰、日间障碍的参照均为``没有或每周不足1次''。*为原文标注的 \(P<0.05\)。Model 1调整年龄、性别、教育、职业、BMI、吸烟、饮酒、午睡、打鼾、自评健康、体力活动、冠心病史、卒中史及糖尿病；Model 2再调整抑郁症状。


\worklines{3}

\end{minipage}\par\vspace{5pt}

\par\addvspace{8pt}\noindent\begin{minipage}{\linewidth}

\section{往年题}

\subsection{单项选择题}

\Question{R24-M12}{Logistic回归系数与优势比}

某二分类Logistic模型没有交互项，自变量 \(X\) 的回归系数为 \(\log 2\)。控制其他变量不变，\(X\) 增加1单位时，下列解释正确的是：


\begin{choices}

\item 事件发生概率一定变为原来的2倍。

\item 事件发生概率一定增加0.5。

\item 事件发生优势变为原来的2倍。

\item 生存时间一定变为原来的2倍。

\item 事件发生优势减少一半。

\end{choices}

\end{minipage}\par\vspace{5pt}

\par\addvspace{8pt}\noindent\begin{minipage}{\linewidth}

\Question{R24-M3}{总体回归系数检验}

Logistic回归中，检验除截距外所有总体偏回归系数是否同时为0，常用下列哪种检验？


\begin{choices}

\item \(t\)检验。

\item \(F\)检验。

\item 似然比检验。

\item \(Z\)检验。

\item Hosmer--Lemeshow检验。

\end{choices}

\end{minipage}\par\vspace{5pt}

\par\addvspace{8pt}\noindent\begin{minipage}{\linewidth}

\subsection{名词解释}

\Question{Z-D3}{优势比（Odds Ratio, OR）}

名词解释：优势比（Odds Ratio, OR）。


\worklines{1}

\end{minipage}\par\vspace{5pt}

\par\addvspace{8pt}\noindent\begin{minipage}{\linewidth}

\subsection{简答题}

\Question{R24-S2}{有序Logistic的阈值与斜率}

如何解释有序分类结局资料的Logistic回归分析中的 \(\alpha_i\) 和 \(X_j\) 的偏回归系数 \(\beta_j\)？


\worklines{3}

\end{minipage}\par\vspace{5pt}

\par\addvspace{8pt}\noindent\begin{minipage}{\linewidth}

\Question{Z-S2}{Logistic偏回归系数的含义}

在Logistic回归中，偏回归系数 \(\beta_j\) 的统计学意义是什么？


\worklines{3}

\end{minipage}\par\vspace{5pt}

\par\addvspace{8pt}\noindent\begin{minipage}{\linewidth}

\Question{Z-S8}{有序结局的共同斜率}

在有序多分类结局的Logistic回归中，偏回归系数 \(\beta_j\) 的解释有何特殊之处？


\worklines{3}

\end{minipage}\par\vspace{5pt}

\par\addvspace{8pt}\noindent\begin{minipage}{\linewidth}

\subsection{计算与综合分析题}

\Question{E20-2}{吸烟与饮酒的交互作用（2020）}

表1 心血管疾病的危险因素分析（Logistic回归）：

\begin{center}\small\setlength{\tabcolsep}{4pt}\renewcommand{\arraystretch}{1.22}\begin{tabular}{@{}lrrrrr@{}}
\toprule
因素 & \(B\) & SE(\(B\)) & CHISQ & \(P\) & OR \\
\midrule

吸烟 vs 不吸烟 & 0.693 & 0.323 & 4.605 & 0.032 & 2.0 \\
饮酒 vs 不饮酒 & 1.099 & 0.551 & 3.975 & 0.046 & 3.0 \\
吸烟\(\times\)饮酒 & 0.405 & 0.186 & 4.752 & 0.029 & 1.5 \\
\bottomrule\end{tabular}\end{center}

利用表1估计：（1）相对于不吸烟、不饮酒者，同时吸烟、饮酒者的发病风险；（2）对于饮酒者，吸烟的发病风险。


\worklines{3}

\end{minipage}\par\vspace{5pt}

\par\addvspace{8pt}\noindent\begin{minipage}{\linewidth}

\Question{E21-2}{吸烟与饮酒的交互作用（2021）}

表1 心血管疾病的危险因素分析（Logistic回归）：

\begin{center}\small\setlength{\tabcolsep}{4pt}\renewcommand{\arraystretch}{1.22}\begin{tabular}{@{}lrrrrr@{}}
\toprule
因素 & \(B\) & SE(\(B\)) & CHISQ & \(P\) & OR \\
\midrule

吸烟 vs 不吸烟 & 0.693 & 0.323 & 4.605 & 0.032 & 2.0 \\
饮酒 vs 不饮酒 & 1.099 & 0.551 & 3.975 & 0.046 & 3.0 \\
吸烟\(\times\)饮酒 & 0.405 & 0.186 & 4.752 & 0.029 & 1.5 \\
\bottomrule\end{tabular}\end{center}

利用表1估计：（1）相对于不吸烟、不饮酒者，同时吸烟、饮酒者的发病风险；（2）对于饮酒者，吸烟的发病风险。


\worklines{3}

\end{minipage}\par\vspace{5pt}

\par\addvspace{8pt}\noindent\begin{minipage}{\linewidth}

\Question{E21-4}{分层分析与年龄混杂}

分析暴露与发病是否有关，判断年龄是否是混杂因素并提供依据。

\begin{center}\small\setlength{\tabcolsep}{4pt}\renewcommand{\arraystretch}{1.22}\begin{tabular}{@{}lrrrrr@{}}
\toprule
结局 & 70--岁 暴露 & 70--岁 未暴露 & 50--岁 暴露 & 50--岁 未暴露 & 合计 \\
\midrule

发病 & 89 & 59 & 80 & 20 & 248 \\
不发病 & 11 & 41 & 120 & 180 & 352 \\
\bottomrule\end{tabular}\end{center}

年龄组标签按原卷``70--岁''\,``50--岁''保留。


\worklines{3}

\end{minipage}\par\vspace{5pt}

\par\addvspace{8pt}\noindent\begin{minipage}{\linewidth}

\Question{R24-C1}{380对配比病例对照研究}

选了380例病例，按照年龄及另外两个未记明的因素选了380例对照，以表的形式报告了 \(\beta\)、\(P\)值、OR及置信区间。

\begin{enumerate}
\def\labelenumi{\arabic{enumi}.}
\tightlist
\item
  本研究的研究设计类型是什么，统计分析方法是什么？（2分）
\item
  该统计分析方法的优点、缺点及用途。（2分）
\item
  模型中没有给出常数项 \(\beta_0\)，需要考虑 \(\beta_0\) 吗？为什么？（2分）
\item
  解释表格中的偏回归系数。（4分）
\end{enumerate}

\textbf{原稿缺失：} 除年龄外的配比因素、全部结果表数值。


\worklines{3}

\end{minipage}\par\vspace{5pt}

\par\addvspace{8pt}\noindent\begin{minipage}{\linewidth}

\subsection{数据分析题}

\Question{E19-3}{低出生体重的影响因素}

利用数据1（原题称 DATA1.sav，目录实际提供 data1.xls 与 data1.csv）分析18岁及以上的母亲所产婴儿是否发生低出生体重的影响因素，并对结果做简单解释。

\begin{center}\small\setlength{\tabcolsep}{4pt}\renewcommand{\arraystretch}{1.22}\begin{tabular}{@{}llll@{}}
\toprule
变量 & 说明 & 变量 & 说明 \\
\midrule

low & 低出生体重：1是、0否 & ftv & 随访次数 \\
age & 产妇年龄（岁） & bwt & 出生体重（g） \\
lwt & 产妇体重（磅） & ui & 应激性：1是、0否 \\
race & 1 white、2 black、3 other & ht & 高血压：1是、0否 \\
smoke & 孕期吸烟：1是、0否 & ptl & 此前早产次数 \\
\bottomrule\end{tabular}\end{center}

数据文件保留了缺失值；请说明分析样本的确定过程。


\worklines{3}

\end{minipage}\par\vspace{5pt}


\clearpage\section{补充练习}

\begin{keyidea}[说明]以下为配合正文新编的补充练习，按题型编排。教学设定的数据不代表真实研究结果；解析见章末。\end{keyidea}

\par\addvspace{8pt}\noindent\begin{minipage}{\linewidth}

\subsection{单项选择题}

\Question{P2-M1}{根据结局与设计选择Logistic模型}

正文的双胞胎配对案例中，结局为随访期间是否新发痴呆（1/0），暴露按抑郁症发生的生命阶段分为7类，并以“终身无抑郁症”为参照。拟控制配对因素，分析暴露与结局的关联，最恰当的方法是：

\begin{choices}

\item 因暴露有7类，采用无序多分类结局的Logistic回归。

\item 按配对组分层的条件Logistic回归，并对暴露设置哑变量。

\item 将暴露编码1—7后，采用有序结局的Logistic回归。

\item 忽略配对关系，把全部个体作为无配对资料直接分析。

\item 以配对编号为连续因变量，作一般线性回归。

\end{choices}

\end{minipage}\par\vspace{5pt}

\par\addvspace{8pt}\noindent\begin{minipage}{\linewidth}

\subsection{名词解释}

\Question{P2-D1}{比例优势假设}

名词解释：比例优势假设（Proportional Odds Assumption）。请结合有序Logistic模型说明。

\worklines{1}

\end{minipage}\par\vspace{5pt}

\par\addvspace{8pt}\noindent\begin{minipage}{\linewidth}

\subsection{简答题}

\Question{P2-S1}{总体检验、增量检验与拟合评价}

正文AMI病例的Logistic分析报告：仅含截距模型的对数似然为\(-122.172\)；含休克、心衰、抢救是否及时三个二分类变量的模型为\(-111.308\)；仅含休克与心衰的模型为\(-115.557\)。各模型使用同一批资料、相同的似然定义。
\begin{enumerate}
\item 写出三个协变量的总体检验原假设、似然比统计量及自由度。
\item 控制休克和心衰后，检验抢救是否及时的关联，应比较哪两个模型？
\item 上述检验显著，能否据此认定模型已通过拟合优度检验？说明理由。
\end{enumerate}
可用\(\chi^2_{0.05,1}=3.841\)、\(\chi^2_{0.05,3}=7.815\)。

\worklines{3}

\end{minipage}\par\vspace{5pt}

\par\addvspace{8pt}\noindent\begin{minipage}{\linewidth}

\subsection{计算与综合分析题}

\Question{P2-C1}{配对四格表中的有效信息}

某1:1配对病例对照研究共80对，只研究一个二分类暴露因素。表中每个数均为\textbf{配对数}。
\begin{center}\begin{tabular}{lrr}\toprule
 & 对照暴露 & 对照未暴露\\\midrule
病例暴露 & 18 & 30\\
病例未暴露 & 10 & 22\\\bottomrule
\end{tabular}\end{center}
\begin{enumerate}
\item 写出暴露系数\(\beta\)的条件似然，求\(\hat\beta\)和配对OR。
\item 为什么两人暴露状态相同的对子不为该暴露系数提供信息？
\item 若忽略配对，把病例、对照分别合并后按普通四格表求OR，结果是多少？应采用哪个结果？
\item 为什么此处不需估计每一配对组的截距？
\end{enumerate}

\worklines{3}

\end{minipage}\par\vspace{5pt}

\par\addvspace{8pt}\noindent\begin{minipage}{\linewidth}

\Question{P2-C2}{无序多分类的概率与参照变换}

正文以泡状细胞癌（\(Y=3\)）为参照，给出的三类肿瘤模型为
\[\begin{aligned}
\log(P_1/P_3)&=-0.9826+0.6281X_1-0.6494X_2,\\
\log(P_2/P_3)&=-0.3461+0.3454X_1-0.6352X_2,
\end{aligned}\]
其中\(Y=1\)为鳞癌，\(Y=2\)为腺癌；\(X_1,X_2\)沿用正文方程中的数值编码。
\begin{enumerate}
\item 对\(X_1=2,X_2=1\)的对象，计算三类结局概率。
\item 保持\(X_2\)不变，\(X_1\)增加1单位时，腺癌相对鳞癌的概率比\(P_2/P_1\)乘以多少？
\item 若分别把两个线性预测子代入二分类公式\(e^g/(1+e^g)\)，能否直接当成\(P_1,P_2\)？
\end{enumerate}

\worklines{3}

\end{minipage}\par\vspace{5pt}

\par\addvspace{8pt}\noindent\begin{minipage}{\linewidth}

\Question{P2-C3}{有序结局的参数个数与解释}

正文的视网膜病变案例将结局记为1无病变、2轻度、3中度、4重度，采用
\[\log\frac{P(Y>j\mid X)}{P(Y\le j\mid X)}
=-\alpha_j+\beta_1X_1+\beta_2X_2+\beta_3X_3+\beta_4X_4.\]
四个自变量为HbA1c（以\%为单位，如7表示7\%）、病程、年龄和性别；均各用一个系数，无交互项。正文报告HbA1c的\(\hat\beta_1=0.63\)、\(SE=0.12\)。
\begin{enumerate}
\item 模型需要估计几个阈值、几个斜率？若改用无序四分类Logistic，共需几个参数？
\item 控制其他变量，HbA1c从7\%升到9\%，处于较高病变等级的累计OR是多少？计算其Wald近似95\%置信区间，并明确比较的两类。
\item 补充设定平行线检验\(P=0.018\)。在\(\alpha=0.05\)下能否继续不加说明地用一个共同OR概括所有切点？
\end{enumerate}

\worklines{3}

\end{minipage}\par\vspace{5pt}

\par\addvspace{8pt}\noindent\begin{minipage}{\linewidth}

\subsection{文献分析题}

\Question{P2-L1}{联合暴露结果表能回答什么}

第8章样题称“某农村队列研究团队发表的一篇文章”，讨论空气污染、BMI与中风。其PM\(_{2.5}\)部分如下，所有OR均以“低暴露且BMI正常”为共同参照，并调整婚姻、教育、收入、饮酒、饮食和体力活动。
\begin{center}\small\begin{tabular}{lll}\toprule
PM\(_{2.5}\) & BMI正常：OR（95\% CI） & BMI偏高：OR（95\% CI）\\\midrule
低暴露 & 1（参照） & 1.25（1.09--1.42）\\
高暴露 & 1.93（1.75--2.12） & 3.08（2.67--3.54）\\\bottomrule
\end{tabular}\end{center}
\begin{enumerate}
\item 仅凭“队列研究团队”和此表，能否确定本篇分析属于前瞻性队列分析？还需核对什么？
\item 若采用含二分类暴露\(E\)、BMI分组\(B\)及\(EB\)交互项的Logistic模型，写出模型，并求BMI偏高者中高暴露相对低暴露的OR。
\item 求乘法交互OR的点估计。仅据三个置信区间均未跨1，能否认为交互作用有统计学意义？
\end{enumerate}

\worklines{3}

\end{minipage}\par\vspace{5pt}


\clearpage\section{习题解析}

\begin{keyidea}[核对方式]按题号核对。缺失选项的补写及新编练习均在对应解析末尾说明。\end{keyidea}

\subsection{作业题}

\Answer{H2-1}{食管癌与饮酒：分层频数资料}

令 \(D=1\) 表示饮酒80+，以0--79为参照；年龄哑变量 \(A_2,A_3,A_4\) 分别表示45--54、55--64、65+，以25--44为参照。按原题频数复算： \[\operatorname{logit}(p)=-2.1515-1.6803D+1.9719A_2+2.4871A_3+2.7409A_4.\]

\begin{longtable}[]{@{}lrrr@{}}
\toprule\noalign{}
变量 & \(\hat\beta\) & SE & OR \\
\midrule\noalign{}
\endhead
\bottomrule\noalign{}
\endlastfoot
饮酒80+ & −1.6803 & 0.1893 & 0.1863 \\
45--54岁 & 1.9719 & 0.3705 & 7.1845 \\
55--64岁 & 2.4871 & 0.3579 & 12.0269 \\
65+岁 & 2.7409 & 0.3627 & 15.5014 \\
\end{longtable}

这些系数的 Wald 检验均 \(P<0.001\)。控制年龄后，饮酒80+组食管癌\textbf{优势}为参照组的0.1863倍；其余系数按同样方式解释。结果方向来自给定频数，不应凭常识改写数据。

总体检验：\(H_0:\beta_D=\beta_{A_2}=\beta_{A_3}=\beta_{A_4}=0\)。用似然比、联合 Wald 或 Score 检验；似然比 \(G^2=2(\ell_{\rm full}-\ell_{\rm null})=197.53\)，\(df=4\)，\(P<0.001\)。逐个检验不能替代严格的联合 Wald 检验。

评价可包括分组 Pearson/偏差、Hosmer--Lemeshow、校准、区分度及 AIC/BIC。原稿报告 HL \(\chi^2=0.740,df=4,P=0.946\)；伪 \(R^2=0.1996\)，逐人展开时 AIC约802。\textbf{分组频数与逐人展开的似然常数不同，AIC和残差偏差不可跨编码直接对比；嵌套模型的似然比差值一致。}

\textbf{校正。} OR不是患病概率之比；BIC应为 \(-2\ell+k\log n\)。原材料 的逐项偏差表与其自身全模型输出不一致，以复算为准。


\Answer{H2-2}{睡眠质量与认知功能：文献分析}

\textbf{（1）原评分口径与原文核对。} 原材料 写``前瞻性随访''；但原文使用2008---2012年第二次体检的15,246名50岁以上参加者，同次评价睡眠与认知。就本篇分析而言，应表述为\textbf{依托队列的横断面分析}，不能把队列平台名称等同于分析设计。

\textbf{（2）} DWRT与MMSE分别为0--10、0--30分的连续评分；DWRT\(<4\)定义记忆障碍，MMSE\(<25\)定义认知功能差。DWRT依次读10词并即时回忆，重复3次，第三次后延迟5分钟回忆，以正确词数计分。原文每词间隔\textbf{1秒}，原评分稿写15秒有误。MMSE涵盖定向、记忆、注意/计算、回忆与语言。

\textbf{（3）} Model 1中入睡\(>60\)分钟、效率\(<65\%\)及每周至少一次日间障碍的 OR分别为1.23、1.26、1.26，有统计学证据；前两类趋势 \(P\) 均为0.04。Model 2中入睡时间与日间障碍关联减弱且区间跨1。效率\(<65\%\)的分类对比原文仍标星，但其\textbf{记忆障碍趋势检验 \(P=0.08\)，不再达到0.05}，不能像原材料那样写趋势仍显著。其他因素均未显示明确关联。应同时说明 OR、区间、趋势与调整层次，避免把优势增幅写成风险增幅。


\subsection{往年题}

\Answer{R24-M12}{Logistic回归系数与优势比}

\textbf{答案：C。}

指数系数 \(e^{\log2}=2\) 是优势比。给定基线概率 \(p\)，变换后的概率为 \(2p/(1+p)\)，通常不等于 \(2p\)。

\textbf{补写说明。} 原稿第12题整题缺失；本题为新编练习，并非恢复的往年原题。


\Answer{R24-M3}{总体回归系数检验}

\textbf{答案：C。}

总体原假设为除截距外所有回归系数同时为0。

\textbf{补写说明。} 原稿A---D保留；E原记为``H****检验''，此处补为Hosmer--Lemeshow检验，不能视为原选项的确证。


\Answer{Z-D3}{优势比（Odds Ratio, OR）}

两个比较条件下事件优势 \(p/(1-p)\) 的比。无交互的Logistic模型中，其他变量固定，\(X_j\)增1单位对应OR=\(e^{\beta_j}\)。


\Answer{R24-S2}{有序Logistic的阈值与斜率}

必须先写模型方向。例如 \[\log\frac{P(Y\le i\mid X)}{P(Y>i\mid X)}=\alpha_i+\sum_j\beta_jX_j.\] \(\alpha_i\) 为 \(X=0\) 时第 \(i\) 个切点的累计对数优势；通常有 \(K-1\) 个阈值。其他变量固定时，\(X_j\) 增加1单位，位于``\(\le i\)''而非``\(>i\)''的累计优势乘以 \(e^{\beta_j}\)，此OR对所有切点相同（比例优势假设）。若写成 \(\alpha_i-X^T\beta\)，正系数对应向高等级移动，解释方向应同步改变。不能笼统说``每升一级的OR''。


\Answer{Z-S2}{Logistic偏回归系数的含义}

其他变量固定时，\(X_j\)增1单位所对应的条件对数优势改变量为 \(\beta_j\)，OR为 \(e^{\beta_j}\)。存在交互时，效应还依赖交互变量水平；不能将OR直接解释为概率比。


\Answer{Z-S8}{有序结局的共同斜率}

比例优势模型对所有累计切点使用共同斜率。先规定 \(\log\{P(Y\le k)/P(Y>k)\}=\alpha_k+X^T\beta\)，则增1单位的累计OR为 \(e^{\beta_j}\)，与切点无关；若使用减号参数化则方向反转。不能解释为固定``上升一级''的概率比。


\Answer{E20-2}{吸烟与饮酒的交互作用（2020）}

令吸烟 \(S\)、饮酒 \(D\) 为0/1，\(\operatorname{logit}p=\beta_0+0.693S+1.099D+0.405SD\)。

（1）\(\mathrm{OR}_{11:00}=e^{0.693+1.099+0.405}\approx2\times3\times1.5=9\)。

（2）在 \(D=1\) 中比较吸烟与不吸烟：\(\mathrm{OR}_{S\mid D=1}=e^{0.693+0.405}\approx3\)。

\textbf{原题用语校正。} 模型能给出优势比，不能据此直接求绝对发病风险；缺少截距时不能计算个体概率。原材料（1）漏掉交互项，写成6，应改为9。


\Answer{E21-2}{吸烟与饮酒的交互作用（2021）}

令吸烟 \(S\)、饮酒 \(D\) 为0/1，\(\operatorname{logit}p=\beta_0+0.693S+1.099D+0.405SD\)。

（1）\(\mathrm{OR}_{11:00}=e^{0.693+1.099+0.405}\approx2\times3\times1.5=9\)。

（2）在 \(D=1\) 中比较吸烟与不吸烟：\(\mathrm{OR}_{S\mid D=1}=e^{0.693+0.405}\approx3\)。

\textbf{原题用语校正。} 模型能给出优势比，不能据此直接求绝对发病风险；缺少截距时不能计算个体概率。原材料（1）漏掉交互项，写成6，应改为9。


\Answer{E21-4}{分层分析与年龄混杂}

粗表为暴露发病169、不发病131；未暴露发病79、不发病221。 \[\mathrm{OR}_{\rm crude}=\frac{169\times221}{131\times79}=3.609.\] 两个年龄层的 OR分别为 \(89\times41/(11\times59)=5.622\)、\(80\times180/(120\times20)=6.000\)，相近。 \[\mathrm{OR}_{\rm MH}=\frac{89\times41/200+80\times180/400}{11\times59/200+120\times20/400}=5.867.\] R复算的连续性校正 MH检验 \(\chi^2=69.567\)，\(P<0.001\)，95\% CI为3.788--9.088；年龄调整后仍有关联。

两层暴露各100、200人，未暴露也各100、200人，因此年龄在暴露组与未暴露组的分布完全相同！在风险比尺度，粗 RR=\([169/300]/[79/300]=2.139\)；按相同年龄权重标准化仍为2.139。粗OR与条件OR不同可由 \textbf{OR非可折叠性} 导致，不能仅凭3.609与5.867差异宣布年龄是混杂因素。

\textbf{答题建议。} 年龄是结局的强预测因素，但此表中与暴露独立，按常见混杂定义没有证据认为年龄混杂了暴露---发病关系。说明效应尺度与非可折叠性；若课程按``粗OR与分层OR差异''判断，应指出该经验标准在本表的局限。


\Answer{R24-C1}{380对配比病例对照研究}

若380对为个体1:1匹配，则为配比病例对照研究，常用条件Logistic回归；若只是频数匹配，应据原设计决定分析，回忆稿不足以绝对确认。

条件Logistic按匹配组条件化，可控制组层面的匹配因素并估计其他因素调整OR；不能估计组内不变的匹配因素效应，信息主要来自组内暴露不一致的组合，且不能直接得到人群绝对发病概率。

各组截距是干扰参数，通过条件似然消去，因此无需另外估计公共 \(\beta_0\)。控制其他协变量及匹配组后，\(\beta_j\) 是 \(X_j\) 增加1单位的对数OR，\(e^{\beta_j}\) 是相应OR；结合95\% CI是否跨1和 \(P\) 值报告。原表缺失，不能给出具体数值解释。


\Answer{E19-3}{低出生体重的影响因素}

\textbf{整理者复算示例，非原卷标准答案。} 先筛选 \(age\ge18\)，按题目以 low 为结局；不将用于定义结局的 bwt 作为解释变量。分类变量设置参照组，使用二项分布、logit连接。原始189条记录中183条满足年龄条件；采用所列协变量与结局的完整记录分析，最终108条、44个事件。未依据出生体重推填缺失 low，也未擅自删除测量值相同的不同编号记录。

模型：\(\operatorname{logit}p=\beta_0+\beta_1age+\beta_2lwt+\)种族哑变量\(+\)吸烟\(+ptl+\)高血压\(+\)应激性\(+ftv\)。

该设定下 lwt 系数−0.0183（\(P=0.048\)）；种族2、3相对1的系数分别1.8377、1.1117（\(P=0.0108,0.0474\)）；吸烟系数1.1530，OR约3.168（\(P=0.039\)）。其余变量在0.05水平无明确证据。缺失比例高，完整记录结果可能受选择偏倚影响，不应当作唯一答案。还应检查共线性、连续变量在logit尺度的函数形式、影响点及拟合。完整R代码和参数表随包提供。


\subsection{补充练习}

\Answer{P2-M1}{根据结局与设计选择Logistic模型}

\textbf{答案：B。} 结局是二分类，且采用个体配对设计，应按配对组建立条件Logistic模型。七分类的是\textbf{暴露变量}，可设6个哑变量；这不使结局变成多分类。模型类别不能只看自变量有几个水平。



\Answer{P2-D1}{比例优势假设}

按正文方向写为
\[\log\frac{P(Y>j\mid X)}{P(Y\le j\mid X)}=-\alpha_j+X^T\beta,
\quad j=1,\ldots,K-1.\]
比例优势假设指同一自变量在所有累计分界点使用同一个斜率。其他条件相同，\(X_i\)增加1单位，对任意切点\(j\)，处于高于该切点的累计优势均乘以\(e^{\beta_i}\)。阈值\(\alpha_j\)可以不同。有的资料将其称为“等比例风险假定”，但这里的英文是\textbf{proportional odds}，应与Cox的比例风险假设区分。



\Answer{P2-S1}{总体检验、增量检验与拟合评价}

\textbf{（1）} \(H_0:\beta_1=\beta_2=\beta_3=0\)，不包括截距。
\[G^2=2[-111.308-(-122.172)]=21.728,\quad df=3.\]
拒绝\(H_0\)，三个协变量整体上改善了相对零模型的拟合。

\textbf{（2）} 比较三变量模型与只含休克、心衰的模型：
\[\Delta G^2=2[-111.308-(-115.557)]=8.498,\quad df=1.\]
拒绝\(H_0:\beta_{\text{及时}}=0\)，提示该变量在控制另外两项后仍有统计学关联。

\textbf{（3）} 不能。与零模型相比有改进、增加一项有价值，均不等于所选模型与实际数据足够一致。绝对拟合评价可用适用条件下的分组Pearson、偏差或Hosmer--Lemeshow检验，并检查残差。偏差\(D=2(\ell_{\rm sat}-\ell_{\rm fit})\)以饱和模型为参照；总体似然比检验以零模型为参照。



\Answer{P2-C1}{配对四格表中的有效信息}

\textbf{（1）} 暴露不一致的对子中，病例暴露的条件概率为\(e^\beta/(1+e^\beta)\)。因此
\[L_c(\beta)=\left(\frac12\right)^{40}
\left(\frac{e^\beta}{1+e^\beta}\right)^{30}
\left(\frac1{1+e^\beta}\right)^{10}.\]
忽略常数，\(\ell_c=30\beta-40\log(1+e^\beta)\)，求导为零得
\[e^{\hat\beta}=\frac{30}{10}=3,\qquad\hat\beta=\log3=1.0986.\]
\textbf{（2）} 暴露一致时组内暴露差为0，该对子贡献为\(1/2\)，与\(\beta\)无关。本题中40对暴露不一致的对子提供暴露效应信息。

\textbf{（3）} 合并后病例暴露48、未暴露32；对照暴露28、未暴露52，所得\(\mathrm{OR}_{\text{忽略配对}}=48\times52/(32\times28)=2.7857\)。本设计应采用利用配对关系的估计3，不能用忽略配对的结果替代。

\textbf{（4）} 给定每对恰有1例病例，组截距在病例与对照的优势比较中消去，条件似然不含这些干扰参数。



\Answer{P2-C2}{无序多分类的概率与参照变换}

\textbf{（1）} \(g_1=-0.3758\)、\(g_2=-0.2905\)。三类概率必须共同归一化：
\[P_3=\frac1{1+e^{g_1}+e^{g_2}},\quad
P_1=\frac{e^{g_1}}{1+e^{g_1}+e^{g_2}},\quad
P_2=\frac{e^{g_2}}{1+e^{g_1}+e^{g_2}}.\]
所得数值约为\(P_1=0.2821\)、\(P_2=0.3072\)、\(P_3=0.4107\)，四舍五入前总和为1。

\textbf{（2）} 两式相减可得\(\log(P_2/P_1)=g_2-g_1\)，故相应概率比乘以
\[\exp(0.3454-0.6281)=\exp(-0.2827)=0.7537.\]
这是\(P_2/P_1\)的变化倍数，不是腺癌自身概率\(P_2\)的变化倍数。

\textbf{（3）} 不能。\(e^{g_1}/(1+e^{g_1})=P_1/(P_1+P_3)\)，它是在只考虑类别1与3条件下的概率，并非三分类总体中的\(P_1\)。正文的多分类模型要求\(P_1+P_2+P_3=1\)。



\Answer{P2-C3}{有序结局的参数个数与解释}

\textbf{（1）} 有序模型有3个阈值和4个共同斜率，共7个参数。无序模型以一类为参照，需3套“截距+4个斜率”，共\(3(1+4)=15\)个参数。

\textbf{（2）} HbA1c增加的是\textbf{2个百分点}，即\(\Delta X_1=2\)。
\[\widehat{\mathrm{OR}}=e^{2\times0.63}=3.5254,\qquad
95\%\,CI=\exp\{2(0.63\pm1.96\times0.12)\}
\approx(2.2025,5.6429).\]
分别比较\(Y>1\)与\(Y\le1\)、\(Y>2\)与\(Y\le2\)、\(Y>3\)与\(Y\le3\)。在比例优势假设下，三种切分对应相同OR。不能写成“恰好升高一级的概率增加到3.53倍”。

\textbf{（3）} 平行线检验的\(H_0\)是各切点斜率相同。\(P=0.018\)提示不满足这一约束，应检查模型设定，考虑允许切点效应不同的模型并据其结果解释。第（2）问是在共同斜率模型成立条件下的计算，不能忽略本项诊断直接作为最终结论。检验值为本题教学设定，并非实测结果。



\Answer{P2-L1}{联合暴露结果表能回答什么}

\textbf{（1）} 不能确定。需看暴露与结局测量的先后、是否排除基线已有中风者、是否随访新发中风，以及本次分析实际使用的资料。队列团队也可能发表基线横断面分析，表中OR不能唯一确定研究设计。

\textbf{（2）} 令\(E=1\)为高暴露、\(B=1\)为BMI偏高，\(Z\)为调整变量：
\[\operatorname{logit}p=\beta_0+\beta_EE+\beta_BB+\beta_{EB}EB+\gamma^TZ.\]
在同一模型及相同调整条件下，BMI偏高者中高暴露相对低暴露的OR为
\[e^{\beta_E+\beta_{EB}}=\frac{3.08}{1.25}=2.464.\]
表中3.08比较的是“双重暴露”与共同参照，不能直接当作BMI偏高层内的空气污染OR。

\textbf{（3）} \(e^{\hat\beta_{EB}}=3.08/(1.93\times1.25)=1.2767\)。它是两层空气污染OR之比。检验交互需检验\(H_0:\beta_{EB}=0\)，使用交互项标准误或含、无交互项的似然比检验。现表未给参数协方差或模型似然，不能求其准确标准误、置信区间或显著性；不能直接相除置信区间端点。
\EndExercises
